% !TEX root = ../thesis.tex
Superconducting circuits are a promising technology towards realizing a large scale quantum computer.
These circuits can protect quantum information against by making the encoded states insensitive to the operators by which noise couples to the circuit.
However, single-mode circuits, such as the transmon or fluxonium, face a choice.
They cannot be simultaneously protected against both relaxation and dephasing when built using only the standard components of a Josephson junction, capacitor, and inductor.
Simultaneous protection requires multi-mode circuits, which are difficult to quantize, simulate, search over, and control.
This thesis addresses each of these in turn.
First, we give a general procedure to go from a circuit diagram to a multimode Hamiltonian, precisely characterizing the residual coordinate freedoms.
% This knob can determine both how the Hamiltonian appears symbolically and how quickly it can be analyzed numerically.
Second, we enumerate all unique circuits built from those standard elements (Josephson junction, capacitor, and inductor).
Using a consistent choice of coordinates, we sort them by the structure of their quantized Hamiltonian and explicitly model all unique three node circuits as novel qubits.
Third, we analyze control of the soft \zp qubit.
We propose the first two qubit gates for the system: a geometric-phase-based CZ and a CNOT that transfers population via conditional splitting on the single qubit gate $\lambda$ system.
Fourth, we consider approximate realizations of a $\cos2\phi$ system.
% , which in its ideal form would break the simultaneous protection restriction of single mode systems.
We derive a symmetry selection rule that determines, from the circuit alone, whether a particular realization forbids first-order flux noise that precludes simultaneous protection.
We apply the test to the \zp, KITE, Harmonium, and rhombus circuits and find that the \zp and Harmonium pass the test, while the KITE and rhombus do not.
Finally, we present GPU-accelerated matrix-free methods for simulating larger multimode circuits.
We close with suggestions for future work.
